\section{Limitations and Conclusion}
\label{sec:conclusion}

\paragraph{Limitations.}
We state four limits explicitly, because they define the regime in which the
claims hold.

\emph{Dense transcription remains out of reach.} As shown in
Sec.~\ref{sec:boundary}, exact character copying (chart axis labels, receipts,
handwritten formulae) does not follow from operator-level modulation, and the
data study we ran---retraining a token-concatenation baseline on identical
data---indicates that the gap is architectural rather than a data artefact. A
practical extension is to retain a small number of pooled visual tokens for
this class of queries only; we deliberately do not include such a variant here,
because it would break the zero-token property that the paper is about.

\emph{The backbone is never adapted.} Keeping the LLM frozen is what buys
bit-exact text retention and cheap training, but it also caps the achievable
accuracy: a $8$B student is no better than a $0.6$B one in our measurements,
which suggests that the adapter interface, not the language capacity, is the
bottleneck. Allowing even a low-rank update of the backbone would likely raise
the ceiling, at the cost of the lossless property.

\emph{Training scale is modest.} Our mixtures contain $21$k and $28$k
short-answer items and a single teacher; we have not studied web-scale
pretraining of the rotation generators. Sample-efficiency curves
(Fig.~\ref{fig:eff}) indicate that distillation makes small mixtures viable,
but we cannot claim behaviour at the scale used by industrial VLMs.

\emph{Two-phase inference.} The tower-free variant runs a bidirectional pass
over patches and then a causal pass over text, which prevents a single fused
kernel from covering both phases. Our latency numbers include both phases on
one GPU, but a production deployment would need a fused implementation to
realise the full benefit.

\paragraph{Conclusion.}
We asked whether a frozen text LLM can be given vision without an external
encoder and without paying for visual tokens in its backbone. \modelname{}
answers with a dual-tier orthogonal rotation architecture---static rotation
with channel scaling in the first $K$ layers to encode the image, then dynamic
block-diagonal rotations generated by multi-frequency or spatial-anchor
queries in the remaining frozen layers---followed by on-policy distillation
over the student's own rollouts. Under identical data, steps and base models,
the resulting models exceed a token-concatenation baseline on visual
perception and hallucination resistance ($44.37$ vs.\ $41.35$ MC2), match it on
the average of perception, hallucination and diagram reasoning ($35.90$ vs.\
$35.87$), beat both single-frequency rotation and LoRA-based built-in encoders
by clear margins, reduce FLOPs by $74.2\%$ and latency by $19.5\%$ at high
resolution, and leave pure-text ability untouched at $70.4\%$ GSM8K and
$71.4\%$ ARC-C. The same recipe transfers to a $0.6$B student. Equally
important, we delimit what the paradigm does not do: operator-level modulation
buys semantic perception, layout reasoning and hallucination resistance, but
not dense character transcription, which we trace to the absence of
per-position key-value addressing. We hope both halves of that statement are
useful---the positive one for building lightweight multimodal models from
frozen LLMs, and the boundary one for deciding when a few visual tokens are
worth their cost.
